\chapter{Dynamic Programming Fundamentals}\label{ch:dp}

Dynamic programming (DP) solves a problem by combining the answers to smaller versions of the same problem, computing each
smaller answer only once. After complexity analysis it is the most important idea in competitive programming, and the one that
takes longest to master, because the hard part is not the code but the choice of \emph{what} to compute. The good news is that
almost every DP is designed by answering the same five questions, and that a handful of classic families (sums of coins, paths in
grids, knapsacks, subsequences, pairs of strings, shapes) cover most of what appears in practice. This chapter develops the method on
those families; Chapter~\ref{ch:dppatterns} adds the patterns built on top of it (bitmasks, intervals, digits, profiles, trees of
states), and Chapter~\ref{ch:dpopt} the optimizations that make a correct but slow DP fast.

\begin{chapterintro}{\gold}{Chapter~\ref{ch:search} (recursion and complete search) and Chapter~\ref{ch:complexity}.}
  \item recognise overlapping subproblems, and see every DP as a computation on a DAG of states;
  \item design a DP by answering five questions: state, transition, base cases, order and answer;
  \item write it top-down or bottom-up, pulling values into a state or pushing them out of it;
  \item solve the classic families: coins and sums, grid paths, knapsacks, increasing subsequences, edit distance and shapes;
  \item reconstruct an optimal solution, shrink the memory to one or two rows, and debug a DP systematically.
\end{chapterintro}

% ======================================================================================================================
\section[What DP is]{What dynamic programming is \hfill \gold}\label{sec:dp-what}

The Fibonacci numbers satisfy $F_0 = 0$, $F_1 = 1$ and $F_n = F_{n-1} + F_{n-2}$. Computed by plain recursion, $F_n$ takes
exponential time, because the recursion solves the same subproblems over and over (Figure~\ref{fig:fibdag}): the call for $F_n$
makes $2F_{n+1} - 1$ calls in total, about $4 \cdot 10^{10}$ for $F_{50}$ (Section~\ref{sec:recursion}). Yet there are only 51
\emph{different} subproblems, $F_0$ to $F_{50}$. Remembering each answer the first time it is computed brings the cost down to 51
additions.

\begin{figure}[htbp]
\centering
\begin{minipage}[c]{9.4cm}\centering
\begin{forest}
  for tree={circle, draw=fgink!85, fill=fgfill, line width=0.7pt, minimum size=6.4mm, inner sep=0pt,
    font=\sffamily\scriptsize, s sep=1.6mm, l sep=4mm, edge={line width=0.7pt, draw=fgline!80!fgink}}
  [$F_5$
    [$F_4$
      [$F_3$, warn [$F_2$, cool [$F_1$] [$F_0$]] [$F_1$]]
      [$F_2$, cool [$F_1$] [$F_0$]]]
    [$F_3$, warn [$F_2$, cool [$F_1$] [$F_0$]] [$F_1$]]]
\end{forest}
\end{minipage}
\hspace{3mm}
\begin{minipage}[c]{5.4cm}\centering
\begin{tikzpicture}
  \foreach \nm/\px/\py in {5/0/0, 3/1.6/0, 1/3.2/0, 4/0.8/-1.2, 2/2.4/-1.2, 0/4.0/-1.2}
    \node[vtx] (f\nm) at (\px, \py) {$F_{\nm}$};
  \foreach \ua/\ub in {5/4, 4/3, 3/2, 2/1, 1/0, 5/3, 3/1, 4/2, 2/0} \draw[ded] (f\ua) -- (f\ub);
  \node[lbl] at (2.0, -2.0) {6 states, 9 dependencies};
\end{tikzpicture}
\end{minipage}
\caption{Left: the recursion tree of $F_5$ has 15 calls, because $F_3$ (orange) is solved twice and $F_2$ (blue) three times.
Right: the same computation as a DAG of distinct states, where an arrow points to a value that is needed; DP evaluates each state
once, in an order that respects the arrows.}
\label{fig:fibdag}
\end{figure}

DP applies when a problem has two properties. \textbf{Overlapping subproblems}: the same smaller problem is needed many times, so
remembering its answer pays off. \textbf{Optimal substructure}: the answer for a problem can be assembled from the answers of its
subproblems; for an optimisation problem, an optimal solution is made of optimal solutions of the parts, and for a counting problem,
the solutions split into disjoint classes that are counted separately and added. When the subproblems do not overlap at all, the
same recursion is plain divide and conquer.

\paragraph{DP is a computation on a DAG.} Draw one vertex per \emph{state} (subproblem) and an arrow from each state to the
states its answer depends on. The arrows never form a cycle, because each state depends only on smaller ones, so the states form a
DAG, and DP evaluates a function on it in topological order. A minimisation DP is a shortest path in this DAG, a maximisation a
longest path, and a counting DP counts paths (Section~\ref{sec:dagpaths}). This view explains the cost, $\Oh{\text{states} +
\text{transitions}}$, and it tells you what to do when the dependencies \emph{do} form cycles: then the problem is a shortest-path
problem on a general graph, solved with Dijkstra or BFS rather than DP.

% ======================================================================================================================
\section[The recipe]{The recipe: five questions \hfill \gold}\label{sec:dp-recipe}

\begin{enumerate}[label=\textbf{\color{cpteal}\arabic*.}, itemsep=2pt]
  \item \textbf{State.} What does $dp[\dots]$ mean, in one precise sentence? (``$dp[s]$ is the number of ordered ways to roll a total
        of $s$.'') Most wrong DPs have a vague state.
  \item \textbf{Transition.} How is a state computed from smaller ones? Look at the \emph{last} decision: the last roll, the last
        coin, the last item taken or skipped.
  \item \textbf{Base cases.} Which smallest states are known directly? ($dp[0] = 1$: the empty sequence.)
  \item \textbf{Order.} In which order must the states be computed so that every state's dependencies are ready? (Increasing $s$.)
  \item \textbf{Answer.} Which state, or combination of states, answers the question? ($dp[n]$.)
\end{enumerate}

The running time is the number of states times the cost of one transition; if that is too slow, shrink one of the two
(Chapter~\ref{ch:dpopt}). The memory is the number of states that must be stored at the same time, which is often far fewer than
all of them (Section~\ref{sec:dp-debug}).

\begin{problem}{Dice Combinations (CSES)}
Count the ways to get the sum $n$ by throwing a die one or more times, where the order of the throws matters, modulo $10^9 + 7$
($n \le 10^6$).
\solution The five answers above: $dp[s] = dp[s-1] + dp[s-2] + \dots + dp[s-6]$, the last throw being $1, 2, \dots$ or 6
(Figure~\ref{fig:dice}). There are $n + 1$ states and six terms per state; the program runs in 0.15\,s for $n = 10^6$.
\end{problem}

\begin{figure}[htbp]
\centering
\begin{tikzpicture}
  \node[blbl, anchor=east] at (-0.2, 0) {$dp[s]$};
  \arrs{D}{0}{0}{1/n, 1/cool, 2/cool, 4/cool, 8/cool, 16/cool, 32/cool, 63/hot, 125/n}
  \idx{D}{9}
  \node[lbl, anchor=east] at (-0.2, -0.53) {$s$};
  \spanabove{D-2}{D-7}{the last throw was 6, 5, 4, 3, 2 or 1}
  \node[notel, anchor=west, text width=5.4cm] at (6.6, -0.05) {$dp[7] = 1 + 2 + 4 + 8 + 16 + 32 = 63$: every sequence with sum
    7 ends with exactly one throw $d$, after a sequence with sum $7 - d$.};
\end{tikzpicture}
\caption{The recipe on Dice Combinations. Each state adds up the six states before it; $dp[0] = 1$ counts the empty sequence.}
\label{fig:dice}
\end{figure}

\begin{code}
void solve() {
    int n;
    cin >> n;
    vi dp(n + 1, 0);  // dp[s]: the number of ordered ways to roll the sum s
    dp[0] = 1;
    rep(s, 1, n + 1) rep(d, 1, 7)  // the last roll is d
        if (d <= s) dp[s] = (dp[s] + dp[s - d]) % MOD;
    print(dp[n]);
}
\end{code}

\subsection{Top-down or bottom-up}

There are two ways to evaluate the same recurrence. \defn{Memoization} (top-down) keeps the recursive definition and adds a cache:
the first call for a state computes and stores its answer, later calls return it. \defn{Tabulation} (bottom-up) fills a table in an
order that respects the dependencies.

\begin{code}
vi memo(n + 1, -1);  // -1: not computed yet
auto fib = [&](auto&& self, int i) -> ll {  // top-down
    if (i <= 1) return i;
    ll& res = memo[i];
    if (res != -1) return res;
    return res = self(self, i - 1) + self(self, i - 2);
};
vi f(n + 1);  // bottom-up
f[1] = 1;
rep(i, 2, n + 1) f[i] = f[i - 1] + f[i - 2];
\end{code}

Memoization is the quickest to write from a recurrence, needs no thought about the order, and computes only the states that are
actually reached, which matters when few of them are. Tabulation has no recursion overhead (often two to five times faster), no
risk of a deep recursion overflowing the stack, and allows the memory tricks of Section~\ref{sec:dp-debug}. A reasonable habit:
think top-down, write bottom-up, and fall back to memoization when the reachable states are sparse or the order is awkward.

\begin{pitfall}[The ``not computed'' marker]
A cache initialised with $-1$ fails silently if $-1$ is a legitimate answer, and one initialised with 0 recomputes every state
whose answer is 0, which can make a memoized DP exponential again. Use a separate \cc{vector<bool> done}, or a value outside the
range of answers.
\end{pitfall}

\subsection{Pulling and pushing}

A transition can be written in two directions (Figure~\ref{fig:pullpush}). In \emph{pull} form, each state reads the states it
depends on: $dp[s] = \sum_d dp[s - d]$. In \emph{push} form, each finished state adds its value into the states that depend on it:
for every $d$, $dp[s + d] \mathrel{+}= dp[s]$. Both visit the same arrows of the DAG and give the same result. Pull is natural when
the predecessors of a state are easy to list; push when the successors are, as in ``from this position you may jump to any of
these'', or when a state has many predecessors but only a few successors.

\begin{figure}[htbp]
\centering
\begin{tikzpicture}
  \node[title, anchor=west] at (0, 1.35) {pull: $dp[s]$ reads};
  \arrs{P}{0}{0}{{}/n, {}/cool, {}/cool, {}/cool, s/hot, {}/n}
  \foreach \cc in {2, 3, 4} \draw[->, line width=0.9pt, draw=fgblue] (P-\cc.north) to[bend left=50] (P-5.north);
  \node[title, anchor=west] at (6.2, 1.35) {push: $dp[s]$ adds itself into};
  \arrs{Q}{6.2}{0}{{}/n, s/hot, {}/cool, {}/cool, {}/cool, {}/n}
  \foreach \cc in {3, 4, 5} \draw[->, line width=0.9pt, draw=fgverm] (Q-2.north) to[bend left=50] (Q-\cc.north);
\end{tikzpicture}
\caption{The same transitions in pull form (left) and push form (right). Pull computes one state from finished ones; push
distributes one finished state forwards.}
\label{fig:pullpush}
\end{figure}

% ======================================================================================================================
\section[Coins and sums]{Coins and sums \hfill \gold}\label{sec:dp-coins}

Coin problems are the cleanest place to see what a state must and must not remember. Given coin values $c_1, \dots, c_k$ (each
usable any number of times) and a target sum $x$, three questions have three different DPs.

\begin{problem}{Minimizing Coins (CSES)}
Find the fewest coins with total $x$, or report that it is impossible ($k \le 100$, $x \le 10^6$).
\solution $dp[s] = 1 + \min_j dp[s - c_j]$ over the coins with $c_j \le s$, with $dp[0] = 0$ and impossible sums marked by a value
larger than any answer (Figure~\ref{fig:mincoins}). Here $10^8$ transitions take 0.21\,s.
\end{problem}

\begin{figure}[htbp]
\centering
\begin{tikzpicture}
  \node[blbl, anchor=east] at (-0.2, 0) {fewest coins};
  \arrs{M}{0}{0}{0/n, 1/n, 2/n, 3/n, 4/cool, 1/n, 2/cool, 1/n, 2/n, 3/n, 2/cool, 3/hot}
  \idx{M}{12}
  \draw[->, line width=0.9pt, draw=fgblue] (M-11.north) to[bend left=70] ($(M-12.north)+(-0.22, 0)$);
  \draw[->, line width=0.9pt, draw=fgblue] (M-7.north) to[bend left=35] (M-12.north);
  \draw[->, line width=0.9pt, draw=fgblue] (M-5.north) to[bend left=30] ($(M-12.north)+(0.22, 0)$);
  \foreach \cl/\tx in {11/{$+1$}, 7/{$+5$}, 5/{$+7$}}
    \node[font=\sffamily\scriptsize\bfseries, text=fgblue, anchor=south east] at ($(M-\cl.north)+(0.02, 0.06)$) {\tx};
  \node[notel, anchor=north west, text width=8.2cm] at (0, -0.85) {Coins $\{1, 5, 7\}$:
    $dp[11] = 1 + \min(dp[10], dp[6], dp[4]) = 1 + \min(2, 2, 4) = 3$, for example $5 + 5 + 1$. Taking the largest coin first
    gives $7 + 1 + 1 + 1 + 1$, five coins.};
\end{tikzpicture}
\caption{Minimizing Coins for every sum up to 11. The last coin is 1, 5 or 7, so $dp[11]$ looks back at three cells.}
\label{fig:mincoins}
\end{figure}

\begin{code}
void solve() {
    int n, x;
    cin >> n >> x;
    vi c(n);
    for (ll& v : c) cin >> v;
    vi dp(x + 1, INF);  // dp[s]: the fewest coins with sum s (INF: impossible)
    dp[0] = 0;
    rep(s, 1, x + 1) for (ll v : c)  // the last coin is v
        if (v <= s && dp[s - v] != INF) chmin(dp[s], dp[s - v] + 1);
    print(dp[x] == INF ? -1 : dp[x]);
}
\end{code}

\begin{pitfall}[Greedy coins]
Taking the largest coin that fits is optimal for ``canonical'' systems such as $\{1, 2, 5, 10, 20, 50\}$, but not in general:
Figure~\ref{fig:mincoins} shows coins $\{1, 5, 7\}$ where greedy uses five coins for 11 and the DP three. Unless the statement
guarantees a canonical system, or you can prove the greedy choice, use the DP.
\end{pitfall}

\subsection{Counting: the order of the loops is the order of the coins}

Counting the ways to form $x$ comes in two flavours. If the order of the coins matters ($2 + 1$ and $1 + 2$ are different), sum over
the \emph{last} coin, exactly as in Dice Combinations. If it does not (Coin Combinations~II), the same recurrence counts every
multiset once for each of its orderings, and the fix is a different state: $dp_j[s]$ is the number of ways to form $s$ using only
the first $j$ coin types. Then $dp_j[s] = dp_{j-1}[s] + dp_j[s - c_j]$ (no coin of type $j$, or at least one), every multiset is
built in one canonical order, types in increasing order, and the two loops are nested the other way round (Figure~\ref{fig:coinorder}).

\begin{figure}[htbp]
\centering
\begin{tikzpicture}
  \node[title, anchor=west] at (-1.9, 1.0) {sums outside, coins inside: ordered};
  \node[blbl, anchor=east] at (-0.2, 0) {$dp[s]$};
  \arrs{O}{0}{0}{1/n, 1/n, 2/n, 3/hot}
  \idx{O}{4}
  \node[notel, anchor=north west, text width=4.6cm] at (-1.9, -0.85) {$s = 3$: $1 + 1 + 1$, $1 + 2$, $2 + 1$.};
  \node[title, anchor=west] at (6.0, 1.0) {coins outside, sums inside: unordered};
  \node[blbl, anchor=east] at (8.6, 0) {coin 1 only};
  \arrs{U}{8.8}{0}{1/n, 1/n, 1/n, 1/n}
  \node[blbl, anchor=east] at (8.6, -0.8) {coins 1 and 2};
  \arrs{V}{8.8}{-0.8}{1/n, 1/n, 2/n, 2/hot}
  \idx{V}{4}
  \node[notel, anchor=north west, text width=5cm] at (6.0, -1.65) {$s = 3$: $\{1, 1, 1\}$ and $\{1, 2\}$.};
\end{tikzpicture}
\caption{Coins $\{1, 2\}$ and sum 3. With the sums in the outer loop, the inner loop chooses the last coin, so every order is
counted. With the coins in the outer loop, coin 2 is added only after all coins 1 are placed, so every multiset is counted once.}
\label{fig:coinorder}
\end{figure}

\begin{problem}{Coin Combinations I (CSES)}
Count the \emph{ordered} ways to form $x$ from the coins, modulo $10^9 + 7$ ($k \le 100$, $x \le 10^6$).
\solution Sums outside, coins inside. Coin Combinations~II, the unordered version, is the worked problem of
Section~\ref{wp:coins}; its code differs only in the order of the two loops.
\end{problem}

\begin{code}
void solve() {
    int n, x;
    cin >> n >> x;
    vi c(n);
    for (ll& v : c) cin >> v;
    vi dp(x + 1, 0);  // dp[s]: ordered sequences of coins with sum s
    dp[0] = 1;
    rep(s, 1, x + 1) for (ll v : c)  // sums outside: the last coin is any v
        if (v <= s && (dp[s] += dp[s - v]) >= MOD) dp[s] -= MOD;
    print(dp[x]);
}
\end{code}

Both values being added are below $10^9 + 7$, so a single subtraction replaces the slower \cc{%} operator in the inner loop: the
program runs in 0.21\,s for $10^8$ transitions.

\begin{tip}
When a counting DP gives too large an answer, the usual cause is that the same object is reached along several paths of the state
DAG. Pin down a \emph{canonical} way to build each object (a fixed order of the parts, the smallest element first, the last change
at a fixed place) and make the state follow it.
\end{tip}

% ======================================================================================================================
\section[Paths in grids]{Paths in grids \hfill \gold}\label{sec:dp-grid}

Monotone paths in a grid, moving only right or down, are the simplest two-dimensional DP: the last step into a cell came from the
cell above or the cell to the left. Counting paths adds the two; a maximum-sum path takes the larger and adds the cell's value; a
minimum-cost path takes the smaller.

\begin{problem}{Grid Paths (CSES)}
An $n \times n$ grid has traps. Count the paths from the top-left to the bottom-right corner that move only right or down and avoid
the traps, modulo $10^9 + 7$ ($n \le 1000$).
\solution $dp[r][c] = dp[r-1][c] + dp[r][c-1]$ for a free cell, 0 for a trap (Figure~\ref{fig:gridpathsdp}); 0.11\,s for
$n = 1000$.
\end{problem}

\begin{figure}[htbp]
\centering
\begin{tikzpicture}
  \tikzset{gc/.style={draw=fgline, line width=0.5pt, minimum size=8mm, inner sep=0pt, font=\ttfamily\small, anchor=center}}
  \foreach \rr/\row in {0/{1,1,1,1}, 1/{1,0,1,2}, 2/{1,1,2,0}, 3/{0,1,3,3}}
    \foreach \val [count=\cc from 0] in \row
      \node[gc] (G\rr\cc) at (0.8*\cc, -0.8*\rr) {\val};
  \foreach \cl in {11, 23, 30} \node[gc, fill=fgink!75] at (G\cl) {};
  \node[gc, hot] at (G33) {3};
  \node[gc, cool] at (G32) {3};
  \draw[->, line width=1.1pt, draw=fgblue] ($(G22.center)+(0.22, -0.12)$) -- ($(G32.center)+(0.22, 0.14)$);
  \draw[->, line width=1.1pt, draw=fgblue] ($(G31.center)+(0.12, -0.22)$) -- ($(G32.center)+(-0.14, -0.22)$);
  \node[notel, anchor=west, text width=6.2cm] at (3.4, -1.2) {Dark cells are traps and hold 0. Each free cell adds the
    cell above and the cell to its left: $3 = 2 + 1$ (blue). The corner holds 3, the number of trap-free monotone paths.};
\end{tikzpicture}
\caption{Grid Paths on the sample grid of the problem.}
\label{fig:gridpathsdp}
\end{figure}

\begin{code}
void solve() {
    int n;
    cin >> n;
    vector<string> g(n);
    for (auto& row : g) cin >> row;
    vvi dp(n, vi(n, 0));  // dp[r][c]: paths from (0, 0) to (r, c)
    rep(r, n) rep(c, n) {
        if (g[r][c] == '*') continue;  // a trap: no path ends here
        if (r == 0 && c == 0) dp[r][c] = 1;
        if (r > 0) dp[r][c] += dp[r - 1][c];  // arrived from above
        if (c > 0) dp[r][c] += dp[r][c - 1];  // arrived from the left
        dp[r][c] %= MOD;
    }
    print(dp[n - 1][n - 1]);
}
\end{code}

\begin{trick}[Two walkers at once]
``Walk from the top-left to the bottom-right and back, collecting values; a cell visited twice counts once'' is not two independent
paths, because the second path depends on what the first collected. Move two walkers forward \emph{simultaneously}, one step each:
after $t$ steps both lie on the anti-diagonal $r + c = t$, so the state $(t, r_1, r_2)$ determines both positions. The walkers are on
the same cell exactly when $r_1 = r_2$, which is where the double count is removed. That is $\Oh{n^3}$ states with four transitions
each.
\end{trick}

% ======================================================================================================================
\section[Knapsack]{Knapsack \hfill \gold}\label{sec:knapsack}

Items have weights $w_i$ and values $v_i$; choose a set of total weight at most $W$ with the largest total value. The last decision
is whether item $i$ is in the set: with $dp[i][c]$ the best value using the first $i$ items and capacity $c$,
\[
  dp[i][c] = \max\bigl(dp[i-1][c],\ dp[i-1][c - w_i] + v_i\bigr) \qquad (\text{the second option only if } w_i \le c),
\]
in $\Oh{nW}$ time. The worked problem Book Shop (Section~\ref{wp:books}) is exactly this, with prices as weights and pages as values.

\paragraph{One row is enough.} Row $i$ depends only on row $i - 1$, so a single array can be updated in place, with one catch:
$dp[c]$ needs the \emph{old} $dp[c - w_i]$. Updating the capacities from high to low keeps that value old; updating them from low to
high would use a value that already includes item $i$, so the item could be taken again and again (Figure~\ref{fig:knapdir}). That is
not a bug but a different problem: the low-to-high loop solves the \emph{unbounded} knapsack, where every item has unlimited copies.

\begin{figure}[htbp]
\centering
\begin{tikzpicture}
  \node[blbl, anchor=east] at (-0.2, 0) {high to low (0/1)};
  \arrs{K}{0}{0}{0/n, 0/n, 3/cool, 3/cool, 3/cool, 3/cool, 3/cool}
  \node[blbl, anchor=east] at (-0.2, -1.5) {low to high (unbounded)};
  \arrs{L}{0}{-1.5}{0/n, 0/n, 3/cool, 3/cool, 6/hot, 6/hot, 9/hot}
  \idx{L}{7}
  \foreach \ua/\ub in {1/3, 2/4, 3/5, 4/6, 5/7} \draw[->, line width=0.8pt, draw=fgline] (K-\ua.north) to[bend left=45] (K-\ub.north);
  \foreach \ua/\ub in {3/5, 4/6, 5/7} \draw[->, line width=0.9pt, draw=fgverm] (L-\ua.north) to[bend left=45] (L-\ub.north);
  \node[notel, anchor=west, text width=5.8cm] at (5.4, -0.7) {One item of weight 2 and value 3 is added to an empty table.
    Going down, every cell reads a cell that does not contain the item yet, so it is used at most once. Going up, cell 4 reads cell
    2, which already contains it (red): the item is taken twice, then three times.};
\end{tikzpicture}
\caption{The direction of the capacity loop is the whole difference between the 0/1 and the unbounded knapsack.}
\label{fig:knapdir}
\end{figure}

\subsection{Counting subsets}

Replacing $\max$ by $+$ counts the subsets instead of optimising over them, and the same direction rule applies.

\begin{problem}{Two Sets II (CSES)}
Count the ways to split $\{1, 2, \dots, n\}$ into two sets with equal sums, modulo $10^9 + 7$ ($n \le 500$).
\solution The total $n(n+1)/2$ must be even; then count the subsets with half the total. Every split is two such subsets (one side
and the other), so the count must be halved, which modulo a prime means multiplying by the inverse of 2. Simpler: leave the number
$n$ out of the DP. Exactly one side of every split lacks $n$, so the subsets of $\{1, \dots, n-1\}$ with half the total are the
splits, each counted once.
\end{problem}

\begin{code}
void solve() {
    int n;
    cin >> n;
    int total = n * (n + 1) / 2;
    if (total % 2) {
        print(0);
        return;
    }
    int half = total / 2;
    vi dp(half + 1, 0);  // dp[s]: subsets of the numbers seen so far with sum s
    dp[0] = 1;
    // leave n out: each split is counted once, by its side without n
    rep(i, 1, n) per(s, i, half + 1) dp[s] = (dp[s] + dp[s - i]) % MOD;
    print(dp[half]);
}
\end{code}

\subsection{Bounded copies}

When item $i$ has $k_i$ copies, treating every copy as a separate item costs $\Oh{W \sum k_i}$. \emph{Binary splitting} replaces the
$k$ copies by bundles of $1, 2, 4, \dots, 2^{t-1}$ copies and a remainder $r = k - (2^t - 1) < 2^t$, each bundle a 0/1 item. Every
count from 0 to $2^t - 1$ is a sum of distinct powers of two, and every count from $2^t - 1$ to $k$ is $r$ plus one of those, so the
bundles reach exactly the counts $0, \dots, k$. The cost drops to $\Oh{W \sum \log k_i}$.

\begin{problem}{Bounded knapsack}
There are $n$ kinds of items; kind $i$ has weight $w_i$, value $v_i$ and $k_i$ copies. Maximise the value within capacity $W$
($n \le 100$, $W \le 10^5$, $k_i \le 1000$).
\end{problem}

\begin{code}
void solve() {
    int n, cap;
    cin >> n >> cap;
    vi dp(cap + 1, 0);  // dp[c]: the best value with total weight at most c
    rep(i, n) {
        ll w, v, k;
        cin >> w >> v >> k;
        // bundles of 1, 2, 4, ... copies and a remainder: every count from
        // 0 to k is a sum of distinct bundles, so 0/1 choices suffice
        for (ll b = 1; k > 0; b *= 2) {
            ll take = min(b, k);
            k -= take;
            per(c, take * w, cap + 1)
                chmax(dp[c], dp[c - take * w] + take * v);
        }
    }
    print(dp[cap]);
}
\end{code}

With $n = 100$ kinds of 1000 copies each and $W = 10^5$ the program makes about $10^8$ updates in 0.15\,s. An exact $\Oh{nW}$
algorithm, which splits the capacities by their remainder modulo $w_i$ and slides a monotone queue along each class, is in
Section~\ref{sec:dpopt-window}.

\usetemplate{Knapsack (0/1, unbounded)}{Ctrl+Alt+T, Shift+Q}{cpknap}

\begin{trick}[When the capacity is huge but the values are small]
With $W$ up to $10^9$ but values summing to at most $10^5$, swap the roles of the two numbers: $dp[v]$ is the smallest weight that
reaches value exactly $v$, and the answer is the largest $v$ with $dp[v] \le W$. Always ask which quantity is small enough to index
the table. For plain subset-sum feasibility, a bitset makes the 0/1 knapsack 64 times faster (Section~\ref{sec:bitset}).
\end{trick}

% ======================================================================================================================
\section[Increasing subsequences]{Increasing subsequences \hfill \gold}\label{sec:lis}

The \defn{longest increasing subsequence} (LIS) of $a_0, \dots, a_{n-1}$ is the longest sequence of indices $i_1 < i_2 < \dots$ with
$a_{i_1} < a_{i_2} < \dots$. The textbook DP takes the last element as the state: $dp[i]$, the length of the longest increasing
subsequence that ends at $a_i$, is $1 + \max\{dp[j] : j < i,\ a_j < a_i\}$, in $\Oh{n^2}$.

\paragraph{The $\Oh{n \log n}$ method.} Keep instead, for every length $k$, the \emph{smallest possible last value} of an increasing
subsequence of length $k + 1$ seen so far; call it $\mathit{tail}[k]$. A smaller last value is never worse, because anything that
extends a subsequence also extends one that ends lower. The tails are strictly increasing (a subsequence of length $k + 2$ contains
one of length $k + 1$ that ends lower), so a new element $x$ either extends the longest subsequence, when it is larger than every
tail, or improves the first tail that is $\ge x$, found by binary search (Figure~\ref{fig:lis}). The number of tails at the end is
the LIS length. The worked problem of Section~\ref{wp:lis} uses exactly this.

\begin{figure}[htbp]
\centering
\begin{tikzpicture}
  \node[blbl, anchor=east] at (-0.2, 0) {$a$};
  \arrs{A}{0}{0}{3/n, 1/good, 4/good, 1/n, 5/good, 9/n, 2/n, 6/good}
  \idx{A}{8}
  \draw[->, line width=0.9pt, draw=fggreen!70!black] (A-8.north) to[bend right=40] (A-5.north);
  \draw[->, line width=0.9pt, draw=fggreen!70!black] (A-5.north) to[bend right=40] (A-3.north);
  \draw[->, line width=0.9pt, draw=fggreen!70!black] (A-3.north) to[bend right=60] (A-2.north);
  \node[lbl, anchor=north west] at (0, -0.85) {predecessor links: $6 \to 5 \to 4 \to 1$};
  \foreach \row/\vals/\txt in {1/{3}/{3: start}, 2/{1}/{1 replaces 3}, 3/{1,4}/{4 extends}, 4/{1,4}/{1 replaces 1},
      5/{1,4,5}/{5 extends}, 6/{1,4,5,9}/{9 extends}, 7/{1,2,5,9}/{2 replaces 4}, 8/{1,2,5,6}/{6 replaces 9}} {
    \arr{T\row}{6.7}{0.75 - \row*0.62}{\vals}
    \node[lbl, anchor=west] at (9.7, 0.75 - \row*0.62) {\txt};
  }
  \node[lbl] at (8.1, 0.65) {tails after each element};
\end{tikzpicture}
\caption{The tails for $3, 1, 4, 1, 5, 9, 2, 6$. The final list has length 4, the LIS length. The tails themselves ($1, 2, 5, 6$) are
not a subsequence; following each element's predecessor, the tail before it at the moment it was placed, gives one: $1, 4, 5, 6$.}
\label{fig:lis}
\end{figure}

\paragraph{Reconstruction.} When $a_i$ is placed at position $k$ of the tails, the element at position $k - 1$ at that moment is
the end of a subsequence of length $k$ that $a_i$ extends; record it as $\mathit{prv}[i]$. Storing indices rather than values in the
tails makes this possible, and following $\mathit{prv}$ from the last tail spells out an LIS backwards.

\begin{problem}{Print a longest increasing subsequence}
Given $n \le 2 \cdot 10^5$ numbers, print the length of a longest strictly increasing subsequence and one such subsequence.
\end{problem}

\begin{code}
void solve() {
    int n;
    cin >> n;
    vi a(n);
    for (ll& x : a) cin >> x;
    // tail[k]: index of the smallest possible last element of an increasing
    // subsequence of length k + 1; prv[i]: the element before a[i] in it
    vi tail, prv(n, -1);
    rep(i, n) {
        auto it = lower_bound(all(tail), a[i],
                              [&](ll j, ll x) { return a[j] < x; });
        int k = it - tail.begin();
        if (k > 0) prv[i] = tail[k - 1];
        if (it == tail.end()) tail.push_back(i);
        else *it = i;
    }
    vi lis;  // walk the predecessors back from the end of the longest one
    for (int i = tail.back(); i != -1; i = prv[i]) lis.push_back(a[i]);
    reverse(all(lis));
    print(sz(lis));
    print(lis);
}
\end{code}

\usetemplate{Longest increasing subsequence}{Ctrl+Alt+T, Shift+L}{cplis}

\begin{trick}[The LIS in disguise]
\begin{itemize}
  \item \textbf{Non-decreasing} subsequences: search with \cc{upper_bound} instead of \cc{lower_bound}, so that equal values extend.
  \item \textbf{Fewest decreasing subsequences} covering an array equals the LIS length (a case of Dilworth's theorem); greedily
        placing each element on the leftmost pile whose top is larger, as in Towers (CSES), builds such a cover.
  \item \textbf{Nested boxes or envelopes}: sort by width increasing and, for equal widths, by height \emph{decreasing}, then take
        the LIS of the heights; the decreasing tie-break stops two boxes of equal width from nesting.
  \item \textbf{Fewest changes} that make an integer array strictly increasing: $n$ minus the longest non-decreasing subsequence
        of $a_i - i$.
\end{itemize}
\end{trick}

% ======================================================================================================================
\section[Two sequences: LCS and edit distance]{Two sequences: LCS and edit distance \hfill \gold}\label{sec:editdist}

With two sequences, the state is a pair of prefixes: $dp[i][j]$ is the answer for the first $i$ elements of $a$ and the first $j$ of
$b$, and the transitions look at the last elements, $a_{i-1}$ and $b_{j-1}$.
\begin{itemize}
  \item \textbf{Longest common subsequence}: if the last elements are equal, match them; otherwise at least one of the two is not
        in the common subsequence:
        \[ dp[i][j] = \begin{cases} dp[i-1][j-1] + 1 & \text{if } a_{i-1} = b_{j-1},\\ \max\bigl(dp[i-1][j],\ dp[i][j-1]\bigr)
           & \text{otherwise.} \end{cases} \]
  \item \textbf{Edit distance} (the fewest insertions, deletions and replacements turning $a$ into $b$): the last operation deletes
        $a_{i-1}$, inserts $b_{j-1}$, or aligns the two, for free if they are equal:
        \[ dp[i][j] = \min\bigl(dp[i-1][j] + 1,\ dp[i][j-1] + 1,\ dp[i-1][j-1] + [a_{i-1} \ne b_{j-1}]\bigr), \]
        with $dp[i][0] = i$ and $dp[0][j] = j$.
\end{itemize}

\begin{figure}[htbp]
\centering
\begin{tikzpicture}
  \tikzset{ec/.style={draw=fgline, line width=0.5pt, minimum size=7mm, inner sep=0pt, font=\ttfamily\small, anchor=center}}
  \foreach \ch [count=\cc] in {M, O, V, I, E}
    \node[font=\ttfamily\small\bfseries, text=fgblue] at (0.7*\cc + 0.7, 0) {\ch};
  \foreach \ch [count=\rr] in {L, O, V, E}
    \node[font=\ttfamily\small\bfseries, text=fgblue] at (0, -0.7*\rr - 0.7) {\ch};
  \foreach \row [count=\rr from 0] in {{0,1,2,3,4,5}, {1,1,2,3,4,5}, {2,2,1,2,3,4}, {3,3,2,1,2,3}, {4,4,3,2,2,2}}
    \foreach \val [count=\cc from 0] in \row
      \node[ec] (E\rr\cc) at (0.7*\cc + 0.7, -0.7*\rr - 0.7) {\val};
  \foreach \cl/\val in {00/0, 11/1, 22/1, 33/1, 34/2, 45/2} \node[ec, cool] at (E\cl) {\val};
  \foreach \ua/\ub in {00/11, 11/22, 22/33, 33/34, 34/45}
    \draw[->, line width=1.1pt, draw=fgverm, shorten <=2.1mm, shorten >=2.1mm] (E\ua.center) -- (E\ub.center);
  \node[notel, anchor=west, text width=5.6cm] at (5.1, -2.2) {Edit distance from \texttt{LOVE} to \texttt{MOVIE} is 2. The
    highlighted path, read from the top-left: replace L by M (diagonal, $+1$), keep O and V (diagonal, $+0$), insert I (right, $+1$),
    keep E. A diagonal step is a match or a replacement, a step right an insertion, a step down a deletion.};
\end{tikzpicture}
\caption{The edit-distance table of \texttt{LOVE} (rows) and \texttt{MOVIE} (columns), and an optimal alignment read from it.}
\label{fig:editdist}
\end{figure}

\begin{problem}{Edit Distance (CSES)}
Find the edit distance of two strings of length at most 5000.
\solution The table has $2.5 \cdot 10^7$ cells; each row depends only on the previous one, so two rows of length $m + 1$ suffice.
That keeps the memory at $\Oh{m}$ instead of 200\,MB of 64-bit numbers, and the program runs in 0.19\,s.
\end{problem}

\begin{code}
void solve() {
    string a, b;
    cin >> a >> b;
    int n = sz(a), m = sz(b);
    // prev[j]: edits turning a[0, i) into b[0, j); only two rows are kept
    vi prev(m + 1), cur(m + 1);
    iota(all(prev), 0);  // from the empty prefix: insert j characters
    rep(i, 1, n + 1) {
        cur[0] = i;  // to the empty prefix: delete i characters
        rep(j, 1, m + 1)  // delete a[i - 1], insert b[j - 1], replace or keep
            cur[j] = min({prev[j] + 1, cur[j - 1] + 1,
                          prev[j - 1] + (a[i - 1] != b[j - 1])});
        swap(prev, cur);
    }
    print(prev[m]);
}
\end{code}

\usetemplate{LCS + edit distance}{Ctrl+Alt+T, [}{cplcs}

\begin{trick}[Reconstruction in linear memory: Hirschberg]
Two rows give the distance but not the alignment, which needs the whole table to walk back. \emph{Hirschberg's algorithm} recovers
it in $\Oh{nm}$ time and $\Oh{n + m}$ memory: run the two-row DP forwards on the first half of $a$ and backwards on the second half,
find the column $j$ where the optimal path crosses the middle row (the one minimising the sum of the two half costs), and recurse
on the two halves. The total work is $nm + nm/2 + nm/4 + \dots < 2nm$.
\end{trick}

% ======================================================================================================================
\section[Choosing the state]{Choosing the state \hfill \gold}\label{sec:dp-state}

The state must contain exactly what the future depends on. A useful test: if two different histories reach the same state but
could continue in different ways, the state is missing information; if the state distinguishes histories that always continue
alike, it is larger than necessary.

\begin{problem}{Array Description (CSES)}
An array of $n$ values in $1..m$ has adjacent values differing by at most 1. Some values are known, the others (shown as 0) are
lost. Count the arrays that match, modulo $10^9 + 7$ ($n \le 10^5$, $m \le 100$).
\solution The index alone is not a valid state, since which values may come next depends on the last value; so the state is the
pair (position, last value). Each layer of states depends only on the previous layer (Figure~\ref{fig:arraydesc}): $\Oh{nm}$ time,
0.11\,s at the limits.
\end{problem}

\begin{figure}[htbp]
\centering
\begin{tikzpicture}
  \tikzset{av/.style={vtx, minimum size=7mm}}
  \foreach \lay/\vals in {0/{0,1,0,0,0}, 1/{1,1,1,0,0}, 2/{0,3,0,0,0}}
    \foreach \val [count=\vv] in \vals
      \node[av] (A\lay\vv) at (2.6*\lay, -0.85*\vv) {\val};
  \foreach \cl in {01, 03, 04, 05, 21, 23, 24, 25} \node[av, dim] at (A\cl) {0};
  \node[av, hot] at (A02) {1};
  \node[av, hot] at (A22) {3};
  \foreach \vv in {1, 2, 3} \node[av, cool] at (A1\vv) {1};
  \foreach \vv in {1, ..., 5} \node[lbl, anchor=east] at ($(A0\vv.west)+(-0.15, 0)$) {value \vv};
  \foreach \ua/\ub in {2/1, 2/2, 2/3} \draw[ded, draw=fgblue] (A0\ua) -- (A1\ub);
  \foreach \ua/\ub in {1/2, 2/2, 3/2} \draw[ded, draw=fgblue] (A1\ua) -- (A2\ub);
  \foreach \lay/\tx in {0/{$a_0 = 2$}, 1/{$a_1 = 0$}, 2/{$a_2 = 2$}} \node[blbl] at (2.6*\lay, -0.05) {\tx};
  \node[notel, anchor=west, text width=5.6cm] at (6.2, -2.55) {Each vertex is a state (position, last value) and shows its
    number of ways. A known value allows a single vertex in its layer (red). The three arrays $2,1,2$, $2,2,2$ and $2,3,2$ are the
    three paths into the last layer.};
\end{tikzpicture}
\caption{Array Description with $n = 3$, $m = 5$ and the array $2, 0, 2$ as a layered DAG of states.}
\label{fig:arraydesc}
\end{figure}

\begin{code}
void solve() {
    int n, m;
    cin >> n >> m;
    vi a(n);
    for (ll& v : a) cin >> v;
    // dp[v]: ways to fill the prefix so far, ending with the value v;
    // dp[0] and dp[m + 1] stay 0 as sentinels
    vi dp(m + 2, 0);
    rep(v, 1, m + 1) dp[v] = a[0] == 0 || a[0] == v;
    rep(i, 1, n) {
        vi nd(m + 2, 0);
        rep(v, 1, m + 1) {
            if (a[i] != 0 && a[i] != v) continue;
            nd[v] = (dp[v - 1] + dp[v] + dp[v + 1]) % MOD;
        }
        dp = nd;
    }
    ll ans = 0;
    rep(v, 1, m + 1) ans = (ans + dp[v]) % MOD;
    print(ans);
}
\end{code}

\paragraph{States that are shapes.} Sometimes the natural subproblem is a smaller copy of the input's \emph{shape}, as when a
rectangle is cut into rectangles.

\begin{problem}{Rectangle Cutting (CSES)}
A rectangle of $a \times b$ unit squares is cut, one straight cut through a whole piece at a time, until every piece is a square. Find
the fewest cuts ($a, b \le 500$).
\solution A first cut splits $w \times h$ into $i \times h$ and $(w - i) \times h$, or $w \times i$ and $w \times (h - i)$, and the
two pieces are then solved independently: $dp[w][h]$ is the minimum over all first cuts of $1 + $ the two pieces' values, with
$dp[w][w] = 0$ (Figure~\ref{fig:rectcut}). Cuts at $i$ and at $w - i$ give the same two pieces, so only cuts up to the middle are
tried, which halves the $\Oh{ab(a + b)}$ work; 0.24\,s for $500 \times 499$.
\end{problem}

\begin{figure}[htbp]
\centering
\begin{tikzpicture}[scale=0.62]
  \fill[fgsky!20] (0, 0) rectangle (3, 3);
  \fill[fggreen!14] (3, 1) rectangle (5, 3);
  \fill[fgorange!16] (3, 0) rectangle (4, 1);
  \fill[fgorange!16] (4, 0) rectangle (5, 1);
  \draw[step=1, draw=fgline!70, line width=0.3pt] (0, 0) grid (5, 3);
  \draw[line width=1pt, draw=fgink] (0, 0) rectangle (5, 3);
  \draw[line width=1.8pt, draw=fgverm, dashed] (3, -0.3) -- (3, 3.3);
  \draw[line width=1.8pt, draw=fgverm, dashed] (3, 1) -- (5.3, 1);
  \draw[line width=1.8pt, draw=fgverm, dashed] (4, -0.3) -- (4, 1);
  \foreach \px/\py/\tx in {3/3.75/1, 5.65/1/2, 4/-0.75/3}
    \node[circle, draw=fgverm, fill=white, line width=0.8pt, inner sep=0.8pt, minimum size=4.2mm,
      font=\sffamily\scriptsize\bfseries, text=fgverm] at (\px, \py) {\tx};
  \node[notel, anchor=west, text width=6.4cm] at (7.0, 1.5) {$3 \times 5$: cut off a $3 \times 3$ square (cut 1); the
    remaining $3 \times 2$ piece loses a $2 \times 2$ square (cut 2); the last $1 \times 2$ strip is cut into two unit squares
    (cut 3). No two cuts suffice, so $dp[3][5] = 3$.};
\end{tikzpicture}
\caption{An optimal way to cut a $3 \times 5$ rectangle into squares.}
\label{fig:rectcut}
\end{figure}

\begin{code}
void solve() {
    int a, b;
    cin >> a >> b;
    // dp[w][h]: the fewest cuts that turn a w x h rectangle into squares
    vvi dp(a + 1, vi(b + 1, INF));
    rep(w, 1, a + 1) rep(h, 1, b + 1) {
        if (w == h) {
            dp[w][h] = 0;
            continue;
        }
        // the first cut, vertical or horizontal; by symmetry, up to the middle
        rep(i, 1, w / 2 + 1) chmin(dp[w][h], dp[i][h] + dp[w - i][h] + 1);
        rep(i, 1, h / 2 + 1) chmin(dp[w][h], dp[w][i] + dp[w][h - i] + 1);
    }
    print(dp[a][b]);
}
\end{code}

% ======================================================================================================================
\section[Reconstruction, memory and debugging]{Reconstruction, memory and debugging \hfill \gold}\label{sec:dp-debug}

\paragraph{Reconstructing a solution.} To output the choices and not only the value, either store, for every state, which transition
achieved the optimum (a parent pointer, as $\mathit{prv}$ in the LIS), or keep the table and re-derive the choice afterwards: walk from
the answer state and, at each step, find a transition whose value plus its cost equals the stored value, exactly as the highlighted path
of Figure~\ref{fig:editdist} was found. The second way costs no extra memory during the DP but needs the full table during the walk.

\paragraph{Shrinking the memory.} If row $i$ of the table depends only on row $i - 1$, keep two rows and swap them (edit distance), or
one row updated in the right direction (knapsack). If a state depends on the last $k$ states only, a ring buffer of $k$ cells is
enough. Each saving gives up reconstruction, unless the choices are stored separately in a compact form, such as one bit per state.

\begin{pitfall}[A checklist for wrong DP answers]
\begin{itemize}
  \item Unreachable states initialised to 0 instead of $-\infty$ or $+\infty$, so that impossible choices look free.
  \item An in-place update in the wrong direction, so that one item is used twice (or a value is read after it changed).
  \item Overflow in a counting DP: reduce modulo after \emph{every} addition and multiplication, and multiply in 64 bits.
  \item A state that is missing information: two histories reach it but need different futures.
  \item Double counting: the same object built along two paths of the state DAG; fix a canonical construction order.
  \item Off-by-one base cases: is $dp[0]$ the empty prefix (usually 1 way, or cost 0), and is the answer $dp[n]$ or $dp[n-1]$?
\end{itemize}
The fastest way to find such bugs is a brute force on tiny inputs and a random test generator (Chapter~\ref{ch:craft}).
\end{pitfall}

\begin{recognize}
Counting ways, a minimum or maximum cost, or ``is it possible'', over a sequence of decisions where the future depends on the past
only through a small summary. Limits such as $n \le 5000$ (for $\Oh{n^2}$), $n \cdot W \le 10^8$, or a grid of $1000 \times 1000$
are typical. A brute-force recursion that keeps being called with the same arguments is a DP waiting to be memoized.
\end{recognize}

\begin{summary}[Chapter summary]
\begin{itemize}
  \item A DP evaluates a function on a DAG of states: the cost is states times transitions, and minimum, maximum and counting DPs
        are shortest paths, longest paths and path counts in that DAG.
  \item Design with five questions: state, transition (look at the last decision), base cases, order, answer.
  \item Memoization is quick to write; tabulation is faster, stack-safe, and allows rolling arrays. Pull and push are the same
        arrows read in two directions.
  \item The loop order of a coin DP decides between ordered and unordered counts; the capacity direction decides between the 0/1
        and the unbounded knapsack; binary splitting handles bounded copies.
  \item The LIS needs only the smallest tail of every length: $\Oh{n \log n}$, and predecessor links reconstruct it.
  \item Pairs of prefixes solve LCS and edit distance; two rows suffice for the value, Hirschberg recovers the alignment.
  \item The state must hold exactly what the future depends on: the last value, the shape that remains.
\end{itemize}
\end{summary}

\begin{practice}
\ppgroup{Counting and optimising sums}
\begin{pplist}
  \item Dice Combinations
  \item Minimizing Coins
  \item Coin Combinations I
  \item Coin Combinations II
  \item Removing Digits
  \item Money Sums
  \item Two Sets II
\end{pplist}
\ppgroup{Grids, knapsacks and sequences}
\begin{pplist}
  \item Grid Paths
  \item Book Shop
  \item Array Description
  \item Edit Distance
  \item Rectangle Cutting
  \item Increasing Subsequence
  \item Towers
\end{pplist}
\end{practice}
